\documentclass[11pt]{article}

\usepackage[margin=1in]{geometry}
\usepackage{graphicx}
\usepackage{booktabs}
\usepackage{hyperref}

\title{Blockchain-Based Environmental Data Integrity for Distributed Monitoring Systems}
\author{KM}
\date{\today}

\begin{document}

\maketitle

\begin{abstract}
Distributed environmental monitoring systems produce time-series records that may pass through multiple sensors, gateways, storage systems, correction workflows, and publication channels. This creates an information-systems integrity problem: users need to determine whether records were modified, deleted, replayed, inserted, corrected without authorization, or disconnected from provenance evidence. This paper evaluates a lightweight blockchain-inspired proof-of-concept for environmental data integrity. Using a bounded OpenAQ API v3 dataset covering Warsaw, Berlin, Paris, and Madrid from 2025-07-01 to 2025-12-31, the study compares four integrity models: conventional storage, audit trail, audit trail with hash-chain linkage, and audit trail with hash-chain linkage plus provenance and permission reconstruction. The experiment uses 112,973 canonical records and 25 controlled tampering scenarios. The resulting threat-coverage matrix shows 20 detected scenarios and 5 expected non-detections, with no missed, partial, or unexpected-alert cases in the implemented scenario set. The contribution is not an environmental air-quality analysis or a production blockchain deployment, but a reproducible information-systems evaluation of how audit trails, hash chains, and provenance reconstruction affect verification coverage for distributed monitoring records.
\end{abstract}

\section{Introduction}

Introduce the research problem, motivation, and contribution.

\section{Related Work}

Summarize relevant literature. Use TODO:CITATION_NEEDED for unsupported claims.

\section{Methodology}

\subsection{Study Design}

The study evaluates data integrity mechanisms for distributed environmental monitoring records. The environmental domain is used as a realistic information-systems substrate: monitoring data are time-series based, multi-source, provenance-sensitive, and often reused outside their original collection context. The study does not evaluate air-quality correctness, pollutant behavior, sensor calibration, or environmental policy effects.

The experiment compares four increasingly expressive integrity models under controlled tampering scenarios. Each model is built from the same canonical dataset, and each applicable threat scenario is injected into the corresponding model artifact. A verifier then examines the tampered artifact and produces alert codes. The evaluation compares verifier alerts against scenario labels and aggregates the results into a threat-coverage matrix.

\subsection{Dataset}

The proof-of-concept uses a bounded extract from OpenAQ API v3. \texttt{TODO:CITATION\_NEEDED} The dataset identifier used in the experiment is \texttt{openaq\_capitals\_2025\_h2}. The selected time window is from 2025-07-01 to 2025-12-31. Four European capital cities were included: Warsaw, Berlin, Paris, and Madrid. For each city, three monitoring locations were selected, producing a total of 12 monitoring locations.

The station selection creates a distributed, multi-source environmental monitoring dataset while keeping the experiment bounded and reproducible. The current experiment treats all selected stations as one combined canonical dataset and one global event stream per integrity model. It does not evaluate cities or stations independently, and it does not perform cross-sensor environmental plausibility checking.

The preprocessing step normalized raw OpenAQ records into canonical measurements. The input contained 116,361 OpenAQ records. After preprocessing, 112,973 canonical records remained and 3,388 records were dropped because the measurement value was missing. The final canonical dataset contained nine measured parameters.

\subsection{Integrity Models}

Model A represents conventional storage only. It stores canonical measurement records without an audit trail, hash-chain linkage, provenance state, or permission state. Its verifier is limited to schema checks and duplicate record identifier checks.

Model B represents an audit trail without hash-chain linkage. Measurements are represented as audit events with deterministic payload hashes and event identifiers. This enables local event-integrity checks, but the event stream is not linked by previous-hash references.

Model C extends Model B with hash-chain linkage. Each event contains a payload hash, a previous hash, and a block hash. This model supports verification of event order and chain continuity.

Model D extends Model C with provenance and permission reconstruction. It includes actor identifiers, key identifiers, permission events, revocation events, correction events, and active key-state reconstruction. This model supports additional checks for actor-related and governance-related integrity failures.

\subsection{Threat Scenarios}

The base scenario set covers five data and stream-integrity threats: value modification, timestamp modification, record deletion, fake record insertion, and replay. These scenarios are applied to Models A-D where the model representation supports the scenario.

Five additional scenarios are specific to Model D: broken provenance, unauthorized correction, revoked actor key usage, missing correction reason, and delayed synchronization. These scenarios require provenance, permission, correction, or synchronization metadata and are therefore outside the implemented scope of Models A-C.

Each implemented scenario uses one deterministic injection per applicable model/threat pair. The current design is a reproducible mechanism-comparison benchmark, not a statistical robustness evaluation over repeated random injections.

\subsection{Evaluation Semantics}

The verifier output is evaluated using scenario labels. A scenario receives the status \texttt{detected} when the expected alert code is present in the verifier output. A scenario receives \texttt{expected\_not\_detected} when the threat is intentionally outside the detection capability of the evaluated model. A \texttt{not\_applicable} cell indicates that the scenario is outside the implemented scope of a model.

The current evaluation reports scenario status counts and a threat-coverage matrix. It does not report precision, recall, F1 score, false-positive rate, or false-negative rate because explicit negative cases were not defined in this experiment.

\section{System Architecture}

\subsection{Proof-of-Concept Scope}

The proof-of-concept is intentionally lightweight. Its purpose is to support the experimental comparison of integrity models rather than to implement a production blockchain platform. The implementation uses Python-based data processing, tabular preprocessing, JSON or JSONL artifacts, SQLite storage where useful, deterministic hashing, controlled tampering generation, model-specific verification, and metrics aggregation.

The system keeps raw environmental measurements outside the integrity layer. Integrity evidence is represented through canonical records, audit events, payload hashes, hash-chain links, provenance metadata, permission events, and verification reports. This separation reflects the central design choice of the paper: the integrity layer should make record manipulation visible without requiring environmental measurements to be stored directly on-chain.

\subsection{Data Flow}

The data flow begins with a bounded OpenAQ extract. Raw records are normalized into canonical measurement records with stable identifiers and standardized fields. The same canonical records are then transformed into four model artifacts. Model A stores conventional measurement records. Model B stores audit events. Model C stores hash-linked audit events. Model D stores hash-linked events enriched with provenance and permission state.

For each applicable threat scenario, a tampering generator creates a modified artifact and a scenario label. The verifier reads the tampered artifact, recalculates relevant hashes, reconstructs stream or permission state where required, and emits alerts. The evaluator compares these alerts against the expected scenario labels and exports per-scenario metrics, aggregate status counts, and the threat-coverage matrix.

\subsection{Hashing Strategy}

The proof-of-concept uses deterministic JSON serialization and SHA-256 hashing for the integrity artifacts. \texttt{TODO:CITATION\_NEEDED} Payload hashes are computed from canonicalized event payloads. In Models C and D, each event also contains a previous-hash reference and a block hash, allowing the verifier to detect disrupted chain continuity.

The hash strategy is used for experimental reproducibility and tamper evidence. It is not presented as a final cryptographic standard selection for production environmental monitoring systems. Any production recommendation about algorithms, signatures, key sizes, or long-term cryptographic validity would require separate standards-based justification. \texttt{TODO:CITATION\_NEEDED}

\subsection{Audit Trail and Provenance Model}

The audit trail represents measurements and governance events as append-oriented events. Measurement events preserve a stable relationship between a canonical record and its hashed payload. Correction and permission events represent changes that should be visible to later verification rather than silently overwriting stored values.

Model D adds provenance and permission reconstruction. Actor and key identifiers are included in events, permission grants and revocations are represented explicitly, and the verifier reconstructs active authorization state over the event stream. This enables checks for unauthorized correction, revoked actor key usage, missing correction reasons, broken provenance, and delayed synchronization.

\subsection{Verification Workflow}

The verification workflow is model-specific. Model A checks schema consistency and duplicate identifiers. Model B checks event payload integrity and event identifiers. Model C additionally checks previous-hash continuity and block hashes. Model D adds provenance, permission, correction, key-state, and synchronization checks.

The verifier does not assess whether a measurement is environmentally plausible or scientifically correct. It only evaluates whether the stored artifact remains internally consistent with the integrity evidence and governance rules represented by the selected model.

\section{Evaluation Results}

\subsection{Evaluation Overview}

This section reports the executed proof-of-concept evaluation for the dataset \texttt{openaq\_capitals\_2025\_h2}. The results should be read as a controlled threat-coverage evaluation of integrity mechanisms, not as an environmental analysis of the monitored locations. The experiment compares four integrity models against a fixed set of tampering scenarios and records whether the verifier produced the expected alert code for each applicable model/threat pair.

Table~\ref{tab:dataset-summary} summarizes the dataset preparation used for the evaluation. The preprocessing stage converted raw OpenAQ records into canonical measurements and excluded records with missing measurement values.

\begin{table}[htbp]
\centering
\caption{Dataset preparation summary for \texttt{openaq\_capitals\_2025\_h2}.}
\label{tab:dataset-summary}
\begin{tabular}{lr}
\toprule
Item & Value \\
\midrule
Input OpenAQ records & 116,361 \\
Canonical records after preprocessing & 112,973 \\
Dropped records & 3,388 \\
Selected monitoring locations & 12 \\
Cities & 4 \\
Monitoring locations per city & 3 \\
Parameters after preprocessing & 9 \\
\bottomrule
\end{tabular}
\end{table}

\subsection{Compared Integrity Models}

Table~\ref{tab:model-summary} summarizes the four compared integrity models. The model sequence is designed to isolate the effect of adding an audit trail, then hash-chain continuity, and finally provenance and permission reconstruction.

\begin{table}[htbp]
\centering
\caption{Integrity models compared in the proof-of-concept.}
\label{tab:model-summary}
\small
\begin{tabular}{lp{0.62\linewidth}}
\toprule
Model & Verification capability \\
\midrule
A & Conventional storage with schema and duplicate record identifier checks. \\
B & Audit trail with event payload and event identifier checks, without hash-chain continuity. \\
C & Audit trail plus hash chain, adding previous-hash and block-hash continuity checks. \\
D & Hash chain plus provenance and permissions, adding actor/key, permission, correction, provenance, and delayed-synchronization checks. \\
\bottomrule
\end{tabular}
\end{table}

\subsection{Executed Scenario Matrix}

The full implemented scenario matrix generated 25 controlled tampering scenarios. For each scenario, the proof-of-concept produced a tampered artifact, a scenario label, a verification output, and a per-scenario evaluation output. The aggregate status counts are shown in Table~\ref{tab:status-counts}.

The run produced 20 detected cases and 5 expected non-detections. No scenarios were marked as missed, partial, or unexpected-alert cases. These counts should be interpreted as scenario-level verification outcomes for the controlled benchmark, not as statistical detection rates.

\begin{table}[htbp]
\centering
\caption{Aggregate scenario status counts.}
\label{tab:status-counts}
\begin{tabular}{lr}
\toprule
Status & Count \\
\midrule
\texttt{detected} & 20 \\
\texttt{expected\_not\_detected} & 5 \\
\texttt{missed} & 0 \\
\texttt{partial} & 0 \\
\texttt{unexpected\_alert} & 0 \\
\bottomrule
\end{tabular}
\end{table}

\subsection{Threat-Coverage Matrix}

Table~\ref{tab:threat-coverage} presents the threat-coverage matrix for Models A-D. The matrix uses three compact status codes: D for detected, END for expected non-detection, and NA for not applicable. A detected cell indicates that the expected alert code was present in the verifier output. An expected non-detection indicates an explicit model limitation rather than an implementation failure. A not-applicable cell indicates that the scenario was outside the implemented scope of the model.

\begin{table}[htbp]
\centering
\caption{Threat-coverage matrix for integrity Models A-D.}
\label{tab:threat-coverage}
\small
\begin{tabular}{p{0.42\linewidth}cccc}
\toprule
Threat & A & B & C & D \\
\midrule
\texttt{value\_modification} & END & D & D & D \\
\texttt{timestamp\_modification} & END & D & D & D \\
\texttt{record\_deletion} & END & END & D & D \\
\texttt{fake\_record\_insertion} & END & D & D & D \\
\texttt{replay} & D & D & D & D \\
\texttt{broken\_provenance} & NA & NA & NA & D \\
\texttt{unauthorized\_correction} & NA & NA & NA & D \\
\texttt{revoked\_actor\_key\_usage} & NA & NA & NA & D \\
\texttt{missing\_correction\_reason} & NA & NA & NA & D \\
\texttt{delayed\_synchronization} & NA & NA & NA & D \\
\bottomrule
\end{tabular}
\end{table}

The model-level coverage counts are summarized in Table~\ref{tab:applicable-coverage}. Not-applicable scenarios are excluded from the applicable-scenario denominator. Expected non-detections are retained as explicit model limitations.

\begin{table}[htbp]
\centering
\caption{Applicable scenario coverage by model.}
\label{tab:applicable-coverage}
\begin{tabular}{lrrrr}
\toprule
Model & Applicable & Detected & Expected not detected & Other \\
\midrule
A & 5 & 1 & 4 & 0 \\
B & 5 & 4 & 1 & 0 \\
C & 5 & 5 & 0 & 0 \\
D & 10 & 10 & 0 & 0 \\
\bottomrule
\end{tabular}
\end{table}

\subsection{Model-Level Interpretation}

Model A detected only the replay scenario in the implemented matrix because replay introduced a duplicate record identifier. It did not contain independent audit or hash evidence for detecting direct value modification, timestamp modification, record deletion, or fake insertion.

Model B detected value modification, timestamp modification, fake record insertion, and replay through audit-event integrity checks. It did not detect record deletion in this implementation because its audit events were not linked into a hash chain.

Model C detected all five base scenarios. The addition of hash-chain linkage made stream-continuity disruption visible, including record deletion.

Model D detected all Model C scenarios and the five additional provenance, permission, correction, revoked-key, and delayed-synchronization scenarios. This result supports the intended role of Model D as the most expressive model in the proof-of-concept, while remaining limited to the controlled scenarios implemented in this benchmark.

\subsection{Reproducibility Artifacts}

The primary machine-readable outputs are the scenario metrics CSV file, the threat-coverage matrix CSV file, the metrics summary JSON file, and the experiment-run manifest. The manifest records artifact paths, file sizes, hashes, scenario counts, and run metadata for the executed scenario matrix.

\section{Discussion}

\subsection{Interpretation}

The experiment shows a progressive increase in threat coverage as additional integrity mechanisms are introduced. Conventional storage provides only limited evidence for post-hoc verification. Adding an audit trail enables detection of local event payload changes, but without hash-chain continuity it remains limited for deletion and stream-order threats. Adding hash-chain linkage makes continuity violations visible. Adding provenance and permission reconstruction extends verification from structural tampering to actor-related and governance-related failures.

The main implication is that blockchain-inspired design should be evaluated at the level of concrete integrity functions rather than as a general claim that blockchain improves environmental monitoring. In the proof-of-concept, the useful mechanisms are audit events, deterministic hashes, previous-hash linkage, explicit permission events, correction metadata, and verifier workflows.

\subsection{Scope of Environmental Interpretation}

The OpenAQ dataset provides a realistic distributed environmental monitoring substrate, but the experiment does not produce environmental-domain conclusions. The results do not compare air quality between Warsaw, Berlin, Paris, and Madrid, do not evaluate pollutant behavior, and do not assess sensor calibration or environmental anomaly detection. The selected monitoring locations support a multi-source dataset for information-systems evaluation.

\subsection{Limitations}

The scenario set is controlled and synthetic. Each scenario uses one deterministic injection per applicable model/threat pair. Therefore, the current evaluation should not be described as a statistical robustness study over many attack positions, stations, timestamps, or random seeds.

The current experiment does not define negative cases. For this reason, precision, recall, F1 score, false-positive rate, and false-negative rate are not reported. The appropriate outputs are scenario status counts, model-level coverage, and the threat-coverage matrix.

The proof-of-concept uses one OpenAQ extract and one global event stream per integrity model. It does not build independent chains per station, city, or gateway. It also does not implement cross-sensor plausibility checking, full intermittent-connectivity simulation, summary-block optimization, or production-grade key management.

\subsection{Future Work}

Future work can extend the benchmark with repeated injections across multiple records, stations, timestamps, and random seeds. A negative-case benchmark would make precision, recall, false-positive rate, and false-negative rate meaningful. Additional work could also evaluate per-station or per-gateway chains, delayed synchronization under offline queues, summary-block verification, and cross-sensor environmental consistency checks as a separate layer from blockchain-inspired integrity verification.

\section{Conclusions}

Summarize findings and outline future work.


\bibliographystyle{plain}
\bibliography{refs/references}

\end{document}
